\chapter{Examples}
\section{Equations and Math Features}\label{sec:equations}

\subsection{Align* and Equation Environment}
\begin{align*}
    a &= b + c \tag{ab} \label{eq:ab} \\
    d &= e + f \dotag \label{eq:de} \\
    g &= h + c \dotag \label{eq:gh}
\end{align*}

Equation \ref{eq:ab} shows the relationship between $a$, $b$, and $c$. Equation \eqref{eq:de} demonstrates another relation. \eqref{eq:gh} is also an equation.\\
Notice that the \texttt{\textbackslash{dotag}} automatically numbers the equation. This numbering will behave exatly like numbering in the equation environment.

\subsection{\textbackslash[ ... \textbackslash]}\label{subsec:integral}
\[
    \int_{a}^{b} f(x) \,dx \dotag \label{eq:integral}
\]

In Section \ref{sec:equations}, we discussed integrals, such as the one in Equation \eqref{eq:integral}.

\section{Symbols and Operators}\label{sec:symbols}

\subsection{Some symbol shorthands}\label{subsec:arrows}

\begin{minted}{latex}
    \begin{gather*}
        \to, \gets, \To, \Gets,\\
        \Nabla, \checkmark, \xmark,  \dotag \\
        \Abs{\br{\abs{ABC}}^{\br{EFG}}}^2_2\\
        \pmat{\1 & 0 \\ 0 & \1} = 
        \bmat{1 &  & \0 \\  & 1 &  \\ \0 &  &  \1} =
        \op{Id}
    \end{gather*}
\end{minted}
compiles to:

    \begin{gather*}
        \to, \gets, \To, \Gets,\\
        \Nabla, \checkmark, \xmark,  \dotag \\
        \Abs{\br{\abs{ABC}}^{\br{EFG}}}^2_2\\
        \pmat{\1 & 0 \\ 0 & \1} = 
        \bmat{1 &  & \0 \\  & 1 &  \\ \0 &  &  \1} =
        \op{Id}
    \end{gather*}

\subsubsection{Underbrace and Overset}
\[
    A = \underbrace{B}_{\text{some description}} + \overset{\text{over}}{C}
\]

\subsubsection{Underset, Overrightarrow and Transpose}
\[
    X^\tp = \underset{\text{under}}{Y} + \overrightarrow{Z}
\]

\subsubsection{Trigonometric and Exponential Functions}
    \begin{align*}
        \sin(\theta) &= \frac{a}{b} \\
        \exp(x) &= e^x
    \end{align*}

\subsubsection{Matrices and Brackets}
    \begin{align*}
        &\text{pmatrix:} \quad \pmat{1 & 2 \\ 3 & 4} \\
        &\text{bmatrix:} \quad \bmat{a & b \\ c & d} \\
        &\text{braket:} \quad \Braket{\Psi | \phi} \\
        &\text{bra-ket:} \quad \bra{\psi} \ket{\phi}
    \end{align*}

\subsubsection{\textbackslash{left(}, \textbackslash{right)}}
we can use the \texttt{\textbackslash br\{\dots\}} command as a shorthand for \texttt{\textbackslash left \dots \textbackslash right}
\[
    \br{\abs{x^{1.5i}}} = \br{\sin\br{\frac{a}{b}}}_{a=0}^{\pi}
\]

\subsubsection{\textbackslash{lim}, \textbackslash{sum}, \textbackslash{cap}, \textbackslash{cup}, \textbackslash{op}, \textbackslash{Op}}
    \begin{align*}
        \Op{lim}_{x \to \infty} f(x) \\
        \sum_{n=1}^{\infty} a_n \\
        A \cap B \\
        C \cup D \\
        \op{Tr}(M) \\
        \Op{argmax}_{x} f(x)
    \end{align*}

\subsubsection{\textbackslash{vec}, \textbackslash{ab}}
    \begin{align*}
        \Nabla{V} &, \abs{a}
    \end{align*}

\section{Theorems and Definitions}

\subsection{Theorems}

\begin{theorem}
    This is a theorem.
\end{theorem}

\begin{definition}
    This is a definition.
\end{definition}

\begin{lemma}
    This is a lemma.
\end{lemma}

\begin{corollary}
    This is a korollar.
\end{corollary}

\subsection{Remarks and Notation}

\begin{notation}
    This is a notation.
\end{notation}

\begin{remark}
    This is a remark.
\end{remark}

We can cite sources: \cite[Additional info]{ftb}
\section{Code Examples}
\subsection{Python Code with \texttt{minted}}
It is possible to do inline syntax highliting 
\begin{minted}{latex}
\mintinline{python}|def foo()|
\end{minted}
or Entire Codeblocks (the lines and frame can both be disabled)
\begin{minted}[frame=single, linenos]{python}
def fibonacci(n):
    if n <= 0:
        return 0
    return fibonacci(n) + fibonacci(n-1)
\end{minted}
\subsection{text with \texttt{minted}}
\begin{minted}[frame=single, linenos]{text}
It is possible to write unformatted text
    like this
    \/§%@ \end{example}
\end{minted}

\subsection{Pseudocode with \texttt{algorithmic}}

\begin{algorithm}[H]
\caption{Euclid}
\begin{algorithmic}[1] % the [1] makes the lines numbered
    \Procedure{Euclid}{$a, b$}\Comment{The greatest common divisor of $a$ and $b$}
        \While{$b \neq 0$}
            \State $r \gets a \bmod b$
            \State $a \gets b$
            \State $b \gets r$
        \EndWhile
        \State \textbf{return} $a$\Comment{The gcd is $a$}
    \EndProcedure
\end{algorithmic}
\end{algorithm}

\section{List Environments}

\subsection{Itemize Environment}

\begin{itemize}
    \item Item 1
    \item Item 2
    \item Item 3
\end{itemize}

\subsection{Enumerate Environment}

\begin{enumerate}
    \item First item
    \item Second item
    \item Third item
\end{enumerate}

\subsection{Customized Enumerate Environment with \texttt{enumitem}}

\begin{enumerate}[(I)]
    \item Customized First item
    \item Customized Second item
\end{enumerate}
\begin{enumerate}[a)]
    \item Customized First item
    \item Customized Second item
\end{enumerate}

\section{Tabular Environment}

\subsection{Basic table with multicolumns and multirows}

\begin{table}[H]
    \centering
    \begin{tabular}{rrrr}
    \multicolumn{4}{l}{Really Long header}\\
    \toprule
    Top & Center & Bottom & Both  \\ 
    \midrule
    \multirow[t]{2}{*}{a}&
    \multirow[c]{2}{*}{b}&
    \multirow[b]{2}{*}{c}&
    d\\
    &&&e \\
    \bottomrule
    \end{tabular}
    \caption{Caption}
    \label{tab:my_label}
\end{table}

\subsection{Multitable example}
\begin{table}[H]
\centering
\begin{subtable}[T]{0.4\textwidth}
\centering
\begin{tabular}{rrr}
\toprule
ang [\si{\deg}] & $m_{lift}$ [\si{g}] & $m_{res}$ [\si{g}]\\
\midrule
20 & 4 & 11 \\
15 & 11 & 17 \\
10 & 22 & 22 \\
5 & 30 & 27 \\
0 & 39 & 29 \\
-5 & 47 & 31 \\
-10 & 55 & 32 \\
-15 & 65 & 32 \\
-20 & 74 & 32 \\
\midrule
$v_{\infty} = 0$ & 70 & 50\\
\bottomrule
\end{tabular}
\caption{Raw data for Measurements of $F_{lift}, F_{res}$}
\end{subtable}
\quad
\begin{subtable}[T]{0.4\textwidth}
\centering
\begin{tabular}{rrr}
\toprule
$N$ & $x$ [mm]& $p_{rel}$ [mm\ce{H2O}]\\
\midrule
0 & 5 & 11.4 \\
1 & 7 & -0.5 \\
2 & 17 & -6.4 \\
3 & 23 & -5.5 \\
4 & 40 & 0.3 \\
5 & 65 & 2.4 \\
6 & 88 & 2.4 \\
7 & 119 & 0.0 \\
8 & 89 & 0.0 \\
9 & 68 & -0.7 \\
10 & 46 & -4.5 \\
11 & 28 & -14.4 \\
12 & 16 & -12.6 \\
13 & 7 & -4.5 \\
\bottomrule
\end{tabular}
\caption{Measurements of the relative pressure}
\end{subtable}
\caption{Raw data}
\end{table}

\section{SIunitx Example}\label{sec:siunitx}

\subsection{Basic Usage}

\begin{itemize}
    \item Length: \SI{2.5(0.050)e-3}{\meter}
    \item Mass: \SI{150}{\gram}
    \item Time: \SI{3.2e-2}{\second}
    \item Temperature: \SI{25}{\degreeCelsius}
\end{itemize}

\subsection{Ranges and Lists}

\begin{itemize}
    \item Range: \SIrange{10}{20}{\degreeCelsius}
    \item List: \SIlist{1;3;5}{\meter\per\second}
\end{itemize}

\subsection{Scientific Notation}

\begin{itemize}
    \item numbers: \num{3.45e5}
    \item values: \SI{3.45e5}{\meter\per\second\squared}
\end{itemize}

\subsection{Angles}

\begin{itemize}
    \item Degrees: \SI{45}{\degree}
    \item Radians: $\SI{2\pi}{\radian}$
\end{itemize}
